% Paper 2, Section 10 — Conclusion (skeleton §10: one paragraph-scale close; the fork
% answer is workload-conditional; the principle as a designer's rule; the bound survives
% as design intuition, not validated theory).
% CLAIMS DISCIPLINE — forbidden here as everywhere: "analog beats digital" (retracted);
% "bound confirmed" / contrapositive claims; uncorrected copy margins; pJ as silicon.

\section{Conclusion}
\label{sec:conclusion}

The question this paper set out to answer --- which datapath should carry a state-space
model on neuromorphic hardware --- turns out not to have a single answer, and the shape
of the non-answer is the contribution. Under a parameter-matched protocol where only the
position of the nonlinearity moves, the three neuromorphic routes reverse their ranking
between a statistical sequence task and a precise-recall task, at matched communication
rates and three seeds on both sides. The regularity behind the reversal is the
datapath-degradation principle: a route is cheap exactly when it degrades the part of the
computation the workload does not depend on. Statistical workloads need an exact graded
output and tolerate a lossy analog state, where the analog route buys a $1.60\times$
proxy-energy reduction for a $+0.160$ bpc (5.3\%) quality cost against a properly
regularized digital baseline. Precise-recall workloads need an exact state and tolerate a
spiked output --- and carrying the state itself in spikes is not a degraded option there
but a non-option, landing at exactly chance at $M{=}32$ and $M{=}64$ while being the
cheapest cell in the table. The energy side, meanwhile, obeys a shape law: the recoverable
proxy energy is set by the readout-to-state ratio $V/H$, and across every shape and budget
tested the recoverable energy and the tolerable quality loss never co-occur --- there is
no quality-matched pJ win in this datapath family at any shape we measured
(\S\ref{sec:hardware-shape}).

For a designer, the principle is an organizing rule, not a forecasting one: used
predictively on two settings it had not seen, it failed both times
(\S\ref{sec:principle-against}), so the honest form of the rule is to \emph{measure} which
side of the datapath --- recurrent state or emitted output --- the workload's loss
actually depends on, and spend the exactness budget there. The
co-design constraints that come with the analog route are equally concrete: the ADC step,
not the send-on-delta threshold, floors the event rate, so event sparsity is purchased
with converter bits rather than a free algorithmic knob, and the threshold does not
transfer across workloads.

The firing-floor bound that motivated this study ends in a different place than intended.
Across a $12\times$ swing in its predicted floor, measured spiking-state activity moved
the other way, and accuracy improved as certified capacity shrank; the bound is not
empirically validated in either direction at this scale. It survives as design intuition
--- the reason to \emph{expect} that a compressed state should not be carried in spikes
--- while the rule itself now rests on the direct measurements of
Section~\ref{sec:copy}: the bound was the reason to expect that result; it is not the
reason to believe it. What would move this from simulation to engineering is measured
energy on real neuromorphic silicon, with the analog state's storage and conversion
priced --- the one number this paper cannot supply and the next one should.
